% Proposed mathematical formulation for the queueing paper.
% The behavioral observable and application-specific contract are to be selected.
% Requires amsmath, amssymb, amsthm, and a proposition environment.
\section{Behavioral constraints for posterior restriction}

Let $D$ denote the observations and let $\pi(d\theta\mid D)$ be a
well-defined posterior on a measurable detailed-model parameter space
$\Theta$. Let $u\in\mathcal U$ specify the workload and operating conditions.
For each $\theta$, define a measurable behavioral observable $B(\theta)$ in
a metric space $(\mathcal B,d)$. The observable includes the conditions in
$\mathcal U$ over which the guarantee is required. Let $\mathcal A$ denote
the assumptions on the workload, queue discipline, initial condition, and
service law under which the models and their observables are defined.

We specify the required behavior using measurable functions
$g_j:\mathcal B\to\mathbb R$, $j=1,\ldots,J$, and define
\[
 F=\{\theta\in\Theta:g_j(B(\theta))\leq0,\ j=1,\ldots,J\}.
\]
We construct a measurable retained set $K\subseteq F$ with
$Z_K:=\pi(K\mid D)>0$ and infer parameters using the restricted posterior
\[
 \pi_K(d\theta\mid D)
 =\frac{\mathbf 1_K(\theta)\pi(d\theta\mid D)}{Z_K}.
\]
Every parameter distributed according to $\pi_K$ then satisfies the
specified behavioral constraints almost surely. The choice of $K$ controls
the range of parameters explored during inference.

\subsection{Contract bounds across the hierarchy}

Let $\hat\theta\in\widehat\Theta$ denote reduced-model parameters, let
$\widehat B(\hat\theta)\in\mathcal B$ denote their measurable behavioral observable,
and let $\hat\pi(d\hat\theta\mid D)$ be the reduced posterior. Suppose
each $g_j$ is $L_j$-Lipschitz with respect to $d$. A contract
$C=(\mathcal A,G_C)$ supplies a jointly measurable upper bound $U_C$ satisfying
\[
 \mathcal A\ \Longrightarrow\
 d\bigl(B(\theta),\widehat B(\hat\theta)\bigr)
 \leq U_C(\theta,\hat\theta)
\]
for every pair on its stated domain. Set $U_C=+\infty$ outside that domain.
For $\epsilon\geq0$, define the certified compatibility relation
\[
 R_C(\theta)=\left\{\hat\theta:
 U_C(\theta,\hat\theta)\leq\epsilon,\quad
 g_j(\widehat B(\hat\theta))+L_j\epsilon\leq0
 \text{ for every }j\right\}.
\]
Its reduced-posterior support and retained level set are
\[
 b_C(\theta)=\int\mathbf1_{R_C(\theta)}(\hat\theta)
                  \hat\pi(d\hat\theta\mid D),\qquad
 K_{C,\gamma}=\{\theta:b_C(\theta)\geq\gamma\},
 \quad 0<\gamma\leq1.
\]

\begin{proposition}
Under $\mathcal A$ and the contract and Lipschitz assumptions above,
$K_{C,\gamma}\subseteq F$. If $\pi(K_{C,\gamma}\mid D)>0$, the posterior
restricted to $K_{C,\gamma}$ satisfies the behavioral constraints almost
surely.
\end{proposition}
\begin{proof}
For $\theta\in K_{C,\gamma}$, positivity of $\gamma$ implies that
$R_C(\theta)$ is nonempty. Choose $\hat\theta\in R_C(\theta)$. For every
$j$, the Lipschitz and contract bounds give
\[
 g_j(B(\theta))
 \leq g_j(\widehat B(\hat\theta))
       +L_j d(B(\theta),\widehat B(\hat\theta))
 \leq g_j(\widehat B(\hat\theta))+L_j\epsilon\leq0.
\]
Thus $\theta\in F$. The posterior statement follows by restriction to
$K_{C,\gamma}$ and normalization by its positive posterior mass.
\end{proof}

For $M\geq1$ reduced-model parameter samples $\hat\theta^{(1)},\ldots,
\hat\theta^{(M)}$, define
\[
 b_{C,M}(\theta)=\frac1M\sum_{m=1}^M
       \mathbf1_{R_C(\theta)}(\hat\theta^{(m)}).
\]
The same inclusion holds for
$K_{C,\gamma,M}=\{\theta:b_{C,M}(\theta)\geq\gamma\}$: membership
provides at least one certified pair. This finite-sample implication uses
the validity of the evaluated contract bounds. Estimation of the population
support $b_C$ is a separate statistical calculation.

A fixed additional screening set $S$ yields
$K=K_{C,\gamma,M}\cap S\subseteq F$, provided $\pi(K\mid D)>0$.
This permits learned screening of proposals while preserving the contract
certificate for every retained parameter. An implementation may check the
certificate for each accepted candidate or prove the inclusion for the
entire region accepted by the learned model.

\subsection{Stochastic queueing observables}

For a stochastic queueing model, let $R_{\theta,u}$ be the response time
of a request under conditions $u$. The observable may be a distribution
functional such as $B(\theta)(u)=\mathbb E[R_{\theta,u}]$, assuming this
expectation exists, or the response-time distribution
$B(\theta)(u,t)=\Pr(R_{\theta,u}\leq t)$. A bound on the former controls
the mean. A probability requirement uses the latter explicitly.

For example, on a finite set of specified workloads, with a deadline $r_u$
and violation tolerance $\delta_u$ for each workload,
the requirement $\Pr(R_{\theta,u}>r_u)\leq\delta_u$ becomes
\[
 g_u(B)=1-B(u,r_u)-\delta_u\leq0.
\]
Under the supremum metric on distribution functions, $g_u$ is
$1$-Lipschitz. A certified discrepancy at most $\epsilon$ and reduced-model
bound $\widehat B(\hat\theta)(u,r_u)\geq1-\delta_u+\epsilon$ therefore
certify the detailed-model requirement. A posterior supported on this
retained set also has posterior-predictive violation probability at most
$\delta_u$, by integrating the conditional violation probabilities.

When a reduced observable is estimated by $\widetilde B(\hat\theta)$,
suppose its error satisfies
$d(\widehat B(\hat\theta),\widetilde B(\hat\theta))\leq e(\hat\theta)$
on a stated coverage event. With the same contract bound on the true
model pair, the computable feasibility condition is
\[
 g_j(\widetilde B(\hat\theta))
       +L_j\bigl(e(\hat\theta)+\epsilon\bigr)\leq0,
 \qquad j=1,\ldots,J.
\]
Statistically estimated contract bounds likewise need valid upper bounds
in the compatibility test. The resulting certificate is conditional on
their simultaneous coverage event. If
that event has probability at least $1-\alpha$, the set-inclusion guarantee
holds with that confidence. Adaptive parameter evaluations require
coverage valid for the selected evaluations, for example a uniform bound
or conditional bounds with an allocated total error probability. A GP
membership probability alone supplies its stated model-based probability;
the certificate requires the corresponding sound or calibrated bound.

\subsection{Evaluation}

We evaluate behavioral error or violation probability under the specified
workloads, the fraction of posterior parameter proposals retained, and
effective samples per unit elapsed time. Elapsed time includes contract
construction, reduced-model inference, learning, and detailed-model
inference. The restricted posterior and its positive normalizing mass
define the sampling target. Convergence of a chosen sampler to that target
requires the assumptions appropriate to its transition kernel.
